\documentclass[a11paper]{article}

\usepackage{karnaugh-map}
\usepackage{subcaption}
\usepackage{tabularx}
\usepackage{libs/titlepage}
\usepackage{libs/document}
\usepackage{booktabs}
\usepackage{multicol}
\usepackage{multirow}
\usepackage{float}
\usepackage{longtable}
\usepackage{varwidth}
\usepackage{graphicx}
\usepackage{siunitx}
\usepackage{pifont}
% \usepackage[toc,page]{appendix}
\usepackage[usenames,dvipsnames]{xcolor}

\title{Rapport d'APP}

\class{Physique des portes logiques}
\classnb{GIF470}

\teacher{Keven Deslandes}

\author{
  \addtolength{\tabcolsep}{-0.4em}
  \begin{tabular}{rcl} % Ajouter des auteurs au besoin
      Benjamin Chausse & -- & CHAB1704 \\
      Shawn Couture    & -- & COUS1912 \\
  \end{tabular}
}

\newcommand{\todo}[1]{\begin{color}{Red}\textbf{TODO:} #1\end{color}}
\newcommand{\note}[1]{\begin{color}{Orange}\textbf{NOTE:} #1\end{color}}
\newcommand{\fixme}[1]{\begin{color}{Fuchsia}\textbf{FIXME:} #1\end{color}}
\newcommand{\question}[1]{\begin{color}{ForestGreen}\textbf{QUESTION:} #1\end{color}}

\newcommand{\quicktab}[4]{
  \begin{table}[H]
    \centering
    \caption{#1}
    \label{tab:#2}
    \begin{tabular}{#3}
      #4
    \end{tabular}
  \end{table}
}

\newcommand{\quickfig}[4]{
\begin{figure}[H]
  \centering
  \includegraphics[width=#3\textwidth]{#4}
  \caption{#1}
  \label{fig:#2}
\end{figure}
}

% Checkboxes
\setlength{\fboxsep}{1pt}
\newcommand{\cbox}{\fbox{\phantom{\ding{51}}}}
\newcommand{\cboxtick}{\fbox{\ding{51}}}%
% self-incrementing Test-ID
\newcounter{tid}
\newcommand{\tid}{\stepcounter{tid}\thetid}

\renewcommand{\frenchtablename}{Tableau}


\begin{document}
\maketitle
\newpage
\tableofcontents
\listoffigures
\listoftables

\newpage
\section{Plan de vérification}

\begin{center}
  \footnotesize
	\begin{longtable}{lp{4cm}p{4cm}p{5cm}l}
		% Headers & Footers {{{
		\caption{Plan de vérification} \label{tab:verif}
		\\

		\toprule
		\multicolumn{3}{l}{Objectif Ciblé} &
		\multicolumn{2}{l}{Test des nouvelles opérations}
		\\

		\midrule
		\#                                 &
		\bfseries Test                     &
		\bfseries Action                   &
		\bfseries Résultat Attendu         &
		\cboxtick
		\\

		\midrule
		\endfirsthead

		\multicolumn{5}{c}%
		{{\itshape \tablename\ \thetable{} -- Continué de la page précédente\ldots}}
		\\

		\midrule
		\#                                 &
		Test                               &
		Action                             &
		Résultat Attendu                   &
		\cboxtick
		\\

		\midrule
		\endhead

		\midrule \multicolumn{5}{r}{{Continué à la prochaine page}}
		\\
		\midrule
		\endfoot

		\bottomrule
		\endlastfoot
		% }}}
		% Tests {{{
		\tid                                       &
    Temps de transition Non-Ou &
		Éxecuter \verb|nor_test.asc| et ourvir les logs &
		Le temps de transition \verb|transition_up| et \verb|transition_down| sont
    en bas de 130ps &
		\cbox \\

	\tid                                         &
    Temps de transition Non-Ou                 &
		Éxecuter \verb|nor_test.asc| et ourvir les logs &
		Le temps de transition \verb|transition_up| et \verb|transition_down| sont
    en bas de 130ps &
		\cbox \\

	  \tid &
    Temps de transition Non-Et &
		Éxecuter \verb|nand_test.asc| et ourvir ses logs &
		Le temps de transition \verb|transition_up| et \verb|transition_down| sont
    en bas de 130ps &
		\cbox \\

		\tid &
    Temps de transition Non-Ou &
		Éxecuter \verb|nor_test.asc| et ourvir ses logs &
		Le temps de transition \verb|transition_up| et \verb|transition_down| sont
    en bas de 130ps &
		\cbox \\

	  \tid &
    Temps de transition Non-Et &
		Éxecuter \verb|nand_test.asc| et ourvir ses logs &
		Le temps de transition \verb|transition_up| et \verb|transition_down| sont
    en bas de 130ps &
		\cbox \\

	  \tid &
    Temps de transition Ou exclusif &
		Éxecuter \verb|xor_test.asc| et ourvir ses logs &
		Le temps de transition \verb|transition_up| et \verb|transition_down| sont
    en bas de 130ps & \\

	  \tid &
    TVm du Non-Et &
		Éxecuter \verb|nand_transfer.asc| et ourvir ses logs &
		Le Vm est proche de VDD/2 (0.9) &
		\cbox \\

	  \tid &
    Vm du Ou exclusif &
		Éxecuter \verb|xor_transfer.asc| et ourvir ses logs &
		Le Vm est proche de VDD/2 (0.9) &
		\cbox \\

    % - Buffer - %

    \tid &
    Temps de propagation &
		Éxecuter \verb|pad_buffer_test.asc| et ourvir ses logs &
		Le temps de propagation est plus rapide que l'autre alternative
    d'inverseurs. &
		\cbox \\

    % - ALU - %

    \tid &
    Largeur de transistors &
		Vérifier les largeurs des transistors &
		Tout les transistors on une longueur de 180n et une largeur égal ou en haut
    de 180n &
		\cbox \\

    % - AND 3 bit - %

    \tid &
    Sortie du ET 3 bits &
		Executer \verb|and3b_test.asc| et analyzer le graphique généré &
		Les sorties respectent la table de vérité d'une porte ET. &
		\cbox \\

    \tid &
    Temps de transition du ET 3 bits &
		Executer \verb|and3b_test_transfer.asc| et ouvrir les logs &
		Le temps de transition entre les entrées et les sorties sont d'au plus
    200ps &
		\cbox \\

    % - ADD 3 bit - %

    \tid &
    Sortie de l'additionneur 3 bits &
		Executer \verb|add3b_test.asc| et analyzer le graphique généré &
		Les sorties respectent la table de vérité d'un additionneur 3 bits. &
		\cbox \\

    \tid &
    Temps de transition du ET 3 bits &
		Executer \verb|add3b_test_transfer.asc| et ouvrir les logs &
		Le temps de transition entre les entrées et les sorties sont d'au plus 200ps &
		\cbox \\


    % }}}
	\end{longtable}
\end{center}

\section{Shématiques diverses}
\label{appdx:schematics}

\subsection{UAL}
\quickfig{Unité d'arithmétique logique (UAL)}{alu}{1}{figures/alu.png}

\newpage
\subsection{Additionneur}
Pour simplifier la schématique et éviter la redondance, une composante
\verb|ADDER_1B| a été créée pour ensuite être utilisée dans l'additionneur
3 Bits.

\quickfig{Additionneur 1 Bit}{add1b}{1}{figures/add1b.png}
\quickfig{Additionneur 3 Bits}{add3b}{.4}{figures/add3b.png}

\newpage
\subsection{Opérateur ET}
\quickfig{ET 3 Bits}{and3b}{1}{figures/and3b.png}


\newpage
\section{Conception du Multiplexeur 2 vers 1}

Afin de simplifier la résolution de sa table de vérité \ref{tab:mux2x1b}, les
fils suivant chaque porte logique on été ajoutés commes étapes intermédiaires.
Ces fils peuvent être vu dans la figure \ref{fig:mux2x1b}.

\begin{minipage}{0.48\textwidth}
\quicktab{Multiplexeur 2 vers 1 Bit}{mux2x1b}{ccccccc}{
  \toprule
  \multirow{2}{*}{$s_{el}$} &
  \multirow{2}{*}{$a$} &
  \multirow{2}{*}{$b$} &
  $T_1$ &
  $T_2$ &
  $T_3$ &
  $o_{ut}$ \\
  & & &
  $\overline{s}$ &
  $\overline{(aT_1)}$ &
  $\overline{(bs)}$ &
  $\overline{(T_2T_3)}$ \\
  \midrule
  0&0&0&1&1&1&0 \\
  0&0&1&1&1&1&0 \\
  0&1&0&1&0&1&1 \\
  0&1&1&1&0&1&1 \\
  1&0&0&0&1&1&0 \\
  1&0&1&0&1&0&1 \\
  1&1&0&0&1&1&0 \\
  1&1&1&0&1&0&1
  \\
  \bottomrule
}
\end{minipage}
\hfill
\begin{minipage}{0.48\textwidth}
\begin{align}
  Y &= \overline{(T_2T_3)} = o_{ut}  \\
  &= \overline{(\overline{[aT_1]}\cdot\overline{[bs]} )} \\
  &= \overline{(\overline{[a\overline{s}]}\cdot\overline{[bs]} )} \\
  &= \overline{\overline{[a\overline{s}]}} + \overline{\overline{[bs]}} \\
  Y &= a\overline{s} + bs \\
  \overline{Y} &= \overline{(a\overline{s} + bs)} \\
  \overline{Y} &= \overline{(a\overline{s})} \cdot \overline{(bs)} \\
  \overline{Y} &= (\overline{a}+\overline{\overline{s}}) \cdot (\overline{b}+\overline{s}) \\
  \overline{Y} &= (\overline{a}+s) \cdot (\overline{b}+\overline{s}) \\
\end{align}
\end{minipage}

\subsection{Dimmensionnement relatif PUN, PDN}

Le dimmensionnement relatif pour PUN et PDN dépend de si les transistors sont
en série ou en parallèle à l'intérieur même des portes logiques. Uniquement
l'inverseur et des portes non et sont utilisé dans le multiplexeur. Hors, dans
l'inverseurs le PUN et le PDN on les mêmes dimmensionnement relatif, tandis que
dans les portes non et, le PDN des NMOS ont un facteur multiplicatif de 2 pour
que leurs temps de transitions soit similaire au PUN à cause que sont PDN
consiste en 2 NMOS en série.

\subsection{Dimensionnement absolue des transistors}

Le dimmensionnement absolue des largeurs de $w_p$ sur $w_n$ à été trouvé avec
l'inverseur. Il s'agissait de trouver des valeurs de marge de bruit (NMH et
NML) pour plusieurs valeurs de mp différentes. Éventuellement, la valeur
d'environ 2.5 donnait que le NMH et le NML était identique et le Vm était de
VDD/2. Ce ratio relatif à été utilisé pour tout les transistors en plus du
dimmensionnement relatif. Les MOSFET PMOS sont donc de base 2.5 fois plus parge
que les MOSFET NMOS.


\newpage
\section{Le tampon de sortie}
\quickfig{Tampon de sortie}{pad_buffer}{.5}{figures/pad_buffer.png}


\begin{align}
  x^n &= \frac{C_L}{C} \\
  x^n &= \frac{1.5\times10^{-12}}{1\times10^{-15}} \\
  x^n &= 1500
\end{align}

On sait que $x=e$, donc:

\begin{equation}
n=\ln(1500) \approx7.31 \\
\end{equation}

On veut un nombre paire pour évité d'inverser le signal. C'est donc soit $8$ ou
$6$. Moins d'inverseur implique un plus petit délai de propagation. Mais plus
d'inverseur = plus de facilité. Cependant, on veut minimiser le délai de
propagation. Il reste donc à calculer la valeur du facteur x pour les deux
valeurs de n et de calculer le plus petit délai de propagation. Donc, pour
$n=6$:

\begin{align}
  x^6 &=1500 \\
  x &=1500^{\frac{1}{6}} \approx3.38 \\
\end{align}

Pour $n=8$:

\begin{align}
  x^8 &= 1500 \\
  x   &= 1500^{\frac{1}{8}} \approx 2.49 \\
\end{align}

Après un essai dans LTspice et des mesures avec des curseurs et des directives,
6 inverseurs donne un délai de propagation moindre que 8 inverseurs. On obtient
$\SI{365}{\pico\s}$ au lieu de $\SI{406}{\pico\s}$ .

\section{Prises de mesures}

\quicktab{Temps de transition des composantes en montée}{ttLH-Sortie}{ll}{
\toprule
Composante & Temps de montée 10\%-90\% ($\SI{}{\pico\s}$) \\
\midrule
AND (\&) & 83 \\
ADD ($+$) & 100 \\ % 100 ns = 100,000 ps
ALU & 95 \\
\bottomrule
}

Les mesures de temps de propagations ont été mesurer à 50\% des signaux au lieu
du 90\% à 10\% pris pour les temps de transitions.

\quicktab{Délais de propagation des composantes}{propagation}{llll}{
\toprule
Composante & Signal & Bas vers Haut ($\SI{}{\pico\s}$) & Haut vers Bas ($\SI{}{\pico\s}$) \\
\midrule
AND (\&) & $a_1 \rightarrow o_1$ & 89 & 101 \\
ADD ($+$) & $a_1 \rightarrow o_1$ & 243 & 237 \\
ADD ($+$) & $c_{in} \rightarrow o_1$ & 555 & 536 \\
\bottomrule
}

\newpage
\appendix
\section{Annexes}

\quickfig{Multiplexeur de 2 entrées de 1 Bits}{mux2x1b}{.5}{figures/mux2x1b.png}
\quickfig{Multiplexeur de 2 entrées de 3 Bits}{mux2x3b}{.3}{figures/mux2x3b.png}

% \quicktab{Additionneur 1 Bit}{add1b}{cccccccc}{
%   \toprule
%   \multirow{2}{*}{$c_{in}$} &
%   \multirow{2}{*}{$a$} &
%   \multirow{2}{*}{$b$} &
%   $T_1$     &
%   $D$       &
%   $T_2$     &
%   $s_{um}$  &
%   $c_{out}$ \\
%   & & &
%   $\overline{(ab)}$ &
%   $a\oplus b$ &
%   $\overline{(Dc_{in})}$ &
%   $D\oplus c_{in}$ &
%   $\overline{(T_1T_2)}$
%   \\
%   0&0&0&1&0&1&0&0 \\
%   0&0&1&1&1&1&1&0 \\
%   0&1&0&1&1&1&1&0 \\
%   0&1&1&0&0&1&0&1 \\
%   1&0&0&1&0&1&1&0 \\
%   1&0&1&1&1&0&0&1 \\
%   1&1&0&1&1&0&0&1 \\
%   1&1&1&0&0&1&1&1 \\
%   \bottomrule
% }


\end{document}
